\subsection{Ausgangsspannung eines Komparators ermitteln}
% Alt: Aufgabe 9

\Aufgabe{
    \label{Aufg9}
    Geben ist der dargestellte Komparator. Die Spannung $U_\mathrm{1}= \mathrm{8,25\,V}$ sei konstant. Wie groß ist die
    Ausgangsspannung für $U_2 = \mathrm{0\,V; 4\,V; 8\,V; 8,2\,V; 8,249\,V; 8,251\,V; 8,3\,V; 10\,V}$?
    \begin{figure}[H]
        \centering
        % Linkes Bild
        \begin{subfigure}[b]{0.45\textwidth}
            \centering
            \input{Tikz/src/FigAusgangsspannungKomparator.tex}\alttext{Das Bild zeigt eine elektrische Schaltung mit einem Operationsverstärker, welcher an eine Versorgungsspannung von Plus und Minus 15 V angeschlossen ist.
                Links vom Operationsverstärker sind zwei Eingangsleitungen: Oben links führt eine Leitung mit Pfeil nach unten, beschriftet mit \(U_1\), welche am positiven Eingang des Operationsverstärkers angeschlossen ist. Darunter eine zweite Leitung ist an den negativen Eingang angeschlossen, beschriftet mit \(U_2\). Zwischen diesen beiden Leitungen ist ein kleiner Pfeil von der oberen zur unteren Leitung mit der Beschriftung \(U_D\).}
        \end{subfigure}
        % Rechtes Bild
        %\begin{subfigure}[b]{0.45\textwidth}
        %\centering
        %\includesvg[width=0.8\textwidth]{Bilder/Aufgaben/FigAusgangsspannungKomparator}
        %\end{subfigure}
        \label{fig:FigAusgangsspannungKomparator}
    \end{figure}
}


\Loesung{
    Ausnutzung der Maschenregel besagt:
    \begin{align*}
        -U_1 + U_2 + U_D & = 0         \\
        U_\mathrm{D}     & = U_1 - U_2
    \end{align*}
    \begin{itemize}
        \item  $U_2 = 0\,V$
              \begin{align*}
                  U_\mathrm{D} & = \mathrm{ 8,25\,V - 0\,V = 8,25\,V} \\
                  U_\mathrm{A} & = \mathrm{13,5\,V}
              \end{align*}
        \item  $U_2 = 4\,V$
              \begin{align*}
                  U_\mathrm{D} & = \mathrm{ 8,25\,V - 4\,V = 4,25\,V} \\
                  U_\mathrm{A} & = \mathrm{13,5\,V}
              \end{align*}
        \item $U_2 = 8\,V$
              \begin{align*}
                  U_\mathrm{D} & = \mathrm{ 8,25\,V - 8\,V = 0,25\,V} \\
                  U_\mathrm{A} & = \mathrm{13,5\,V}
              \end{align*}
        \item $U_2 = 8,2\,V$
              \begin{align*}
                  U_\mathrm{D} & = \mathrm{ 8,25\,V - 8,2\,V = 0,05\,V} \\
                  U_\mathrm{A} & = \mathrm{8\,V}
              \end{align*}
        \item $U_2 = 8,249\,V$
              \begin{align*}
                  U_\mathrm{D} & = \mathrm{ 8,25\,V - 8,249\,V = 1\,mV} \\
                  U_\mathrm{A} & = \mathrm{4\,V}
              \end{align*}
        \item $U_2 = 8,251\,V$
              \begin{align*}
                  U_\mathrm{D} & = \mathrm{ 8,25\,V - 8,251\,V = -1\,mV} \\
                  U_\mathrm{A} & = \mathrm{-4\,V}
              \end{align*}
        \item $U_2 = 8,3\,V$
              \begin{align*}
                  U_\mathrm{D} & = \mathrm{ 8,25\,V - 8,3\,V = -0,05\,V} \\
                  U_\mathrm{A} & = \mathrm{-8\,V}
              \end{align*}
        \item $U_2 = 10\,V$
              \begin{align*}
                  U_\mathrm{D} & = \mathrm{ 8,25\,V - 10\,V = -1,75\,V} \\
                  U_\mathrm{A} & = \mathrm{-13,5\,V}
              \end{align*}
    \end{itemize}
    Siehe: Abschnitt \ref{fig: Verschaltung Verstaerker} und Tabelle \ref{tab:Grundschaltungen}
}